\section{Стислі теоретичні відомості}

Під час виконання лабораторної роботи, а саме проєктування системи було використано 
такі фундаментальні концепції об'єктно-орієнтованого програмування:
\begin{itemize}
    \item \textbf{Успадкування (\texttt{extends}):} 
    Механізм, що дозволяє створювати нові класи на основі вже існуючих, успадковуючи 
    їхні поля та методи, забезпечуючи повторне використання коду. В Java підтримується 
    лише поодиноке успадкування для класів.
    \item \textbf{Абстрактні класи (\texttt{abstract}):} 
    Класи, призначені для створення спільної основи. Вони можуть містити поля, конструктори, 
    звичайні методи та абстрактні методи без реалізації, які обов'язково мають бути перевизначені
     в дочірніх класах.
    \item \textbf{Інтерфейси (\texttt{interface} / \texttt{implements}):} 
    Визначають певний набір методів (контракт поведінки) без реалізації. Клас може реалізовувати 
    декілька інтерфейсів одночасно, що компенсує відсутність множинного успадкування класів.
    \item \textbf{Поліморфізм:} 
    Здатність програми працювати з об'єктами різних дочірніх класів через спільний базовий 
    тип (суперклас або інтерфейс), забезпечуючи динамічне зв'язування методів на етапі виконання.
\end{itemize}
